% Poster Trekie — versione LaTeX (A4, 2 colonne)
% Template di riferimento: latex/model.tex (rho class) — qui usiamo KOMA-Script (scrartcl)

\documentclass[10pt,a4paper]{scrartcl}

\usepackage[T1]{fontenc}
\usepackage[utf8]{inputenc}
\usepackage[italian]{babel}
\usepackage{geometry}
\usepackage{graphicx}
\usepackage{enumitem}
\usepackage{lmodern}
\usepackage{microtype}
\usepackage{xcolor}
\usepackage{multicol}
\usepackage[hidelinks]{hyperref}

\geometry{top=12mm,bottom=12mm,left=12mm,right=12mm}

\setlength{\parindent}{0pt}
\setlength{\parskip}{0.35em}
\setlength{\columnsep}{10mm}

\setlist[itemize]{leftmargin=*,topsep=0pt,parsep=0pt,partopsep=0pt,itemsep=3.75pt}

\definecolor{accent}{RGB}{25,55,95}

\pagestyle{empty}

% Spaziatura e stile titoli più adatti a un poster
\RedeclareSectionCommand[beforeskip=0.9ex plus 0.2ex minus 0.1ex, afterskip=0.35ex]{section}
\RedeclareSectionCommand[beforeskip=0.55ex plus 0.2ex minus 0.1ex, afterskip=0.25ex]{subsection}

\setkomafont{section}{\large\bfseries\color{accent}}
\setkomafont{subsection}{\normalsize\bfseries\color{accent}}

\newcommand{\PosterSection}[1]{%
  \section*{#1}%
  \vspace{0.4em}%
  {\color{accent!35}\hrule height 0.4pt}%
  \vspace{0.25em}%
}

\begin{document}

\begin{multicols*}{2}

\begin{center}
  {\Huge\sffamily\bfseries\color{accent} Trekie}\\[-0.05em]
  {\large\sffamily Un’app di navigazione dell’assicurazione sanitaria}\\[0.35em]
  
  \includegraphics[width=\columnwidth,height=40mm,keepaspectratio]{../img/logo.png}\\[.5cm]
  \textbf{Idea chiave:} una funzionalità aggiuntiva che descrive, passo dopo passo,
  come orientarsi nel sistema sanitario pubblico e privato. 
\end{center}

\PosterSection{CPS}
\subsection*{Cliente}
Studenti internazionali, nomadi digitali e viaggiatori di breve durata che vivono o soggiornano in un paese straniero.

\subsection*{Problema}
Questi utenti si trovano ad affrontare termini assicurativi complessi e sistemi sanitari sconosciuti, con conseguente confusione durante le emergenze e un alto rischio di rifiuto delle richieste di rimborso o costi inattesi.

\subsection*{Soluzione}
Un’app di navigazione che traduce le polizze in linguaggio semplice e fornisce guide passo-passo personalizzate al cliente per accedere alle cure e per ottenere e inviare le richieste di rimborso.

\PosterSection{Roadmap}
\begin{itemize}
  \item \textbf{Anno 0 — Sviluppo app:} un’app di navigazione dell’assicurazione sanitaria.
  \item \textbf{Mese 6 — Acquisizione utenti:} obiettivo 10.000 utenti; focus su marketing, costruzione della community e partnership.
  \item \textbf{Anno 1 — Vendita dati:} dati aggregati + anonimizzati + dati comportamentali (non identificabili).
  \item \textbf{Anno 2 — Vendita licenza} a università/aziende.
\end{itemize}

\PosterSection{Piano di lavoro}
\subsection*{Principio}
Vincere prima un pilot piccolo e misurabile in collaborazione con un dipartimento, poi usare quel caso come riferimento per approcciare le università più grandi.

\subsection*{Selezione target}
\begin{itemize}
  \item Preferire un’istituzione privata/più piccola (es. NABA/IED/Domus) o un singolo dipartimento/program di un grande ateneo (meno stakeholder, approvazioni più rapide).
  \item Definire uno scope iniziale stretto: “navigazione sanitaria in una data città, fornire passi per emergenze e spiegazione SSN vs assicurazione privata, per fasi”
\end{itemize}

\PosterSection{Punti di forza (Unique Selling Proposition)}
\subsection*{Per gli utenti dell’app}
\begin{itemize}
  \item Localizzazione
  \item Guida passo-passo durante le emergenze
  \item Personalizzazione della procedura, in base a sintomi, gravità, tipo di assicurazione personale e azienda assicuratrice.
\end{itemize}
  
\subsection*{Per le assicurazioni}
\begin{itemize}
  \item Insight di alta qualità
  \item Dati di utilizzo - domande più frequenti
  \item Analisi per la segmentazione di una nuova offerta di mercato
\end{itemize}

\subsection*{Per le università}
\begin{itemize}
  \item Riduzione del carico operativo di Welcome Office e segreterie grazie a guide personalizzate e chiare.
  \item Meno escalation e richieste ripetute: studenti più autonomi nel navigare assicurazione e sistema sanitario locale.
  \item Miglior esperienza e retention degli studenti internazionali tramite supporto multilingue nel momento del bisogno.
\end{itemize}
  
\subsection*{Funzionalità incluse}
\begin{itemize}
  \item Guida informativa passo-passo alla raccolta documenti
  \item Catalogo di centri medici convenzionati e dei MMG che parlano inglese
  \item Invio e gestione delle richieste di rimborso
\end{itemize}


\PosterSection{Modello di ricavi}
Alle università viene proposto un pagamento di una licenza bastata sul numero di studenti utilizzatori. Gli utenti finali hanno la possibilità di pagare un \textit{premium} per ottenere funzionalità aggiuntive. 

\subsection*{Licenza studenti}
\begin{itemize}
  \item \textbf{Prezzo:} 10\,€ per studente
  \item \textbf{Quantità stimata:} circa 2.000 studenti internazionali per università
  \item \textbf{Obiettivo secondo anno:} 2-3 università
  \item \textbf{Obiettivo quarto anno:} 7-25 università
\end{itemize}



\PosterSection{Interviste}
\begin{itemize}
  \item \textbf{10} interviste tenute con gli studenti internazionali di Trento
  \item \textbf{20} risposte ad un survey mirato sui bisogni degli studenti
  \item Validazione con il \textit{Welcome Office} dell'università di Trento (con Dott. Nicola Battelli)
  \item Validazione con il Research and Development Manager di Dedalus Italia
\end{itemize}








\PosterSection{Attrattività del mercato}
\begin{itemize}
  \item L’Italia è un mercato in forte crescita e poco servito.
  \item 2,05M studenti internazionali totali nel 2024/25 — massimo storico.
  \item +10\% crescita annua degli studenti internazionali dal 2022.
  \item \#3 destinazione UE con crescita più rapida
\end{itemize}


\PosterSection{Opportunità per Trekie}
Il mercato italiano degli studenti internazionali è il più in crescita in Europa, con \textbf{zero concorrenza} nella navigazione assicurativa post-acquisto e multilingue. Trekie entra in uno spazio incontestato, con punti critici misurabili per le università usati come leva di vendita.

\PosterSection{Adeguatezza al mercato (market fit)}
\subsection*{TAM}
\textbf{Mercato delle piattaforme di navigazione sanitaria:} 12,64B\,€

\subsection*{SAM}
\textbf{632M\,€} — viaggiatori/studenti internazionali + expat + nuovi arrivati in Italia, 5\% del TAM

\subsection*{SOM}
\textbf{6,32M\,€} — 1\% del SAM

\PosterSection{Business Model Canvas}
\begin{center}
  \includegraphics[width=\columnwidth,height=40mm,keepaspectratio]{../img/nice graph.jpg}
\end{center}

\subsection*{Segmenti di clientela}
Creiamo valore per:
\begin{itemize}
  \item Grandi aziende che richiedono frequenti viaggi e che organizzano l’assicurazione per i propri dipendenti.
  \item Grandi organizzazioni che gestiscono e assicurano un numero elevato di persone.
  \item Individui che viaggiano/vivono all’estero per lavoro o studio e sono assicurati dalla propria azienda/organizzazione.
\end{itemize}
I clienti più importanti:
\begin{itemize}
  \item Aziende multinazionali che negoziano costantemente con altre aziende all’estero.
  \item Aziende di viaggio che operano su più paesi (compagnie aeree, autobus, spedizioni, ecc.).
  \item Università con alta mobilità internazionale.
\end{itemize}

\subsection*{Proposta di valore}
Riduciamo la complessità organizzativa nella gestione delle interazioni tra le assicurazioni sanitarie e i loro studenti, lavoratori e collaboratori. Il valore si basa su un design human-centric che offre informazioni chiave e non distorte nel momento del bisogno, grazie a una prioritizzazione fondata su conoscenza empirica continua degli utenti target.

Aiutiamo a mitigare informazioni poco chiare, eterogenee e verbose nel settore delle assicurazioni sanitarie.

Il prodotto core è un’app che fornisce informazioni empiricamente/statisticamente rilevanti per diversi tipi di utenti (dagli studenti ai dipendenti full-time). Per il settore sanitario/assicurativo privato, offriamo insight aggregati sui bisogni e sulle sfide degli utenti.

Soddisfiamo il bisogno delle università di avere studenti consapevoli dell’assicurazione e in grado di navigare i servizi sanitari locali in modo fluido e autonomo; aiutiamo le aziende a ridurre la complessità legata alle frizioni dei sistemi sanitari che impattano i dipendenti; e aiutiamo le compagnie assicurative a generare una domanda più alta e più mirata per i loro prodotti.

\subsection*{Canali}
Raggiungiamo i clienti tramite:
\begin{itemize}
  \item Contatto diretto con le aziende.
  \item Pubblicità/articoli su riviste per imprenditori e dirigenti.
\end{itemize}

\subsection*{Relazioni con i clienti}
\begin{itemize}
  \item Per i dirigenti: report/dashboard di controllo su eventuali incidenti, avanzamento e assistenza clienti.
  \item Per i dipendenti: l’app Trekie con supporto in emergenza e per le richieste assicurative, oltre al supporto clienti.
\end{itemize}

\subsection*{Flussi di ricavi}
\begin{itemize}
  \item \textbf{Abbonamenti corporate:} 30\,€ al mese per dipendente.
  \item \textbf{Abbonamenti università:} 10\,€ al mese per studente.
  \item \textbf{Contratti enterprise:} ricavi ricorrenti da accordi B2B di lungo periodo.
\end{itemize}

\subsection*{Risorse chiave}
\begin{itemize}
  \item Ecosistema Trekie (app)
  \item Talento tecnico e commerciale
  \item Infrastruttura di supporto
  \item Asset digitali (cloud e algoritmi di tracciamento)
\end{itemize}

\subsection*{Attività chiave}
\begin{itemize}
  \item Interviste e analisi utenti
  \item Design CHI continuo
  \item Sviluppo app
  \item Networking e costruzione della presenza
  \item Vendita e networking in presenza (eventi mirati/meeting specifici)
\end{itemize}
Le relazioni con i clienti sono il nucleo del modello e vengono mantenute dai fondatori e da un team di esperti commerciali. I ricavi derivano dall’interazione continua con i clienti.

\subsection*{Partner chiave}
Gruppi di consulenza per le relazioni con i clienti, compagnie assicurative e (se possibile) incubatori di startup.

I fornitori chiave includono aziende tech che forniscono framework e librerie per lo sviluppo; i gruppi di consulenza forniscono conoscenza e insight vitali. I partner sviluppano la tecnologia che abilita funzionalità distintive legate all’interazione uomo-computer (CHI) e quindi vantaggi competitivi. Inoltre, offrono conoscenza continua che aiuta a costruire una base di conoscenza solida e coerente.

\subsection*{Struttura dei costi}
\begin{itemize}
  \item \textbf{R\&D piattaforma:} sviluppo software, manutenzione e cybersecurity.
  \item \textbf{Acquisizione \& marketing:} stipendi del team vendite e pubblicità B2B mirata.
  \item \textbf{Operazioni:} stipendi staff di supporto 24/7 e costi infrastruttura cloud.
  \item \textbf{Legale:} compliance assicurativa internazionale e normative privacy.
\end{itemize}

\end{multicols*}

\end{document}
